\chapter{CH50 Complement Activity}

\section*{What Is This Test?}
The CH50 test measures the ability of a patient's serum to activate the classical complement pathway and lyse antibody-coated red cells. It provides a functional assessment of the classical pathway rather than measuring one complement protein alone.

\section*{Alternative Names}
\begin{itemize}
  \item CH50
  \item Total Hemolytic Complement
  \item Classical Pathway Complement Activity
\end{itemize}
\section*{Why This Test Is Ordered}
CH50 supports evaluation of recurrent bacterial infections, suspected inherited complement deficiency, systemic lupus erythematosus, immune-complex disease, vasculitis, and monitoring of complement consumption. C3 and C4 are often measured at the same time.

\section*{How This Test Is Performed}
A blood sample is collected, separated promptly, and kept at the required temperature because complement activity is unstable. The assay determines the serum dilution that produces a defined degree of lysis.

\section*{Preparation, Risks, and Limitations}
No fasting is required. Venipuncture is the usual risk. Delayed separation, warming, or poor transport can produce a falsely low result. A low CH50 may reflect consumption, inherited deficiency, or specimen deterioration and requires correlation with C3, C4, and clinical findings.

\section*{Understanding Results}
Absent CH50 suggests deficiency or severe depletion of one or more classical-pathway components. A low value with low C3 or C4 may indicate active immune-complex disease. A normal CH50 does not exclude alternative-pathway disorders or every complement-mediated disease.

\section*{References}
Clinical laboratory guidance on complement testing in immune-complex disease and complement deficiency.

\clearpage
\section*{Understanding Results and Follow-up}
The laboratory report should identify the specimen, method, units, reference interval or decision limit, and whether the result is positive, negative, abnormal, or inconclusive. Reference intervals vary with age, sex, pregnancy, specimen type, assay, and laboratory; a result outside a range is not automatically a diagnosis, and a result within a range does not always exclude disease. Discuss unexpected results with the ordering clinician, who can decide whether repeat or confirmatory testing, treatment, or referral is appropriate. Do not change medicines or delay urgent care based on an isolated result.


\section*{Regulatory Status and Authorization Date}
This chapter describes a test category, not a specific commercial product. FDA, CDSCO, EU IVDR, and MHRA approval or registration dates therefore cannot be assigned without the assay manufacturer, product name, instrument, and intended use. For laboratory-developed tests, an individual agency approval date may not exist. See the Preface for the interpretation of regulatory status.

\clearpage
